%% lampiran/C-kode-implementasi.tex
%% Lampiran C — Tautan Kode Implementasi dan Reproducibility

\chapter{Tautan Kode Implementasi dan Reproducibility}
\label{lamp:kode}

Seluruh kode implementasi, konfigurasi eksperimen, dan \emph{checkpoint}
model yang digunakan dalam disertasi ini tersedia secara publik untuk
mendukung \emph{reproducibility}.

\section{Repositori}
\label{lamp:kode_repo}

\begin{itemize}
  \item Repositori utama: \url{https://github.com/USERNAME/disertasi-mambaxlstm}
  \item Lisensi: MIT (kode), CC-BY-4.0 (dokumentasi)
  \item \emph{Commit hash} yang sesuai dengan hasil pada Bab~V:
        \texttt{XXXXXXX}
\end{itemize}

\section{Cara Menjalankan}
\label{lamp:kode_run}

% Snippet command-line untuk reproduksi eksperimen E1, E2, dst.
\begin{verbatim}
# Setup environment
conda env create -f environment.yml
conda activate disertasi

# Reproduksi eksperimen E1 (perbandingan SOTA pada XJTU-SY)
python scripts/run_experiment.py --config configs/E1_xjtusy.yaml --seed 42

# Reproduksi seluruh eksperimen (5 seed × 5 eksperimen)
bash scripts/run_all.sh
\end{verbatim}

\section{Daftar Versi Pustaka}
\label{lamp:kode_versi}

% Salin dari requirements.txt / environment.yml
\dots
